\documentclass[10pt,a4paper]{amsart}
\usepackage[margin=19mm]{geometry}
\usepackage{amsmath,amssymb,mathtools}
\usepackage[scr=rsfs,cal=euler]{mathalfa}
\usepackage{xfrac}
\usepackage[unicode,hidelinks,bookmarks=false]{hyperref}
\newcommand{\ie}{i.e.\ }
\newcommand{\cf}{cf.\ }
\newcommand{\resp}{respectively\ }
\newcommand{\Subsection}{\subsection}
\providecommand{\nobreakdash}{\nobreak\hbox{-}}
\setlength{\emergencystretch}{2em}
\multlinegap0pt
\usepackage{amssymb}
\usepackage{enumerate}
\usepackage[shortlabels]{enumitem}
\newtheorem{lemma}{Lemma}
\newcommand{\R}{{\mathbb{R}}}
\newcommand{\dd}{\text{\rm d}}             % a straight d for differentials
\newcommand{\e}{\varepsilon}
\title{Feynman--Kac formula for the heat equation with a one-center point interaction in $d=3$.}
\author{Makoto Nakashima}
\date{Source: arXiv:2606.11677v1 (CC BY 4.0)}
\begin{document}
\maketitle

\noindent\emph{Source and licence.} Feynman--Kac formula for the heat equation with a one-center point interaction in $d=3$.. Version arXiv:2606.11677v1 (CC BY 4.0), distributed by arXiv under the Creative Commons Attribution 4.0 International licence, http://creativecommons.org/licenses/by/4.0/, as recorded in the arXiv licence metadata for this exact version and shown on the abstract page https://arxiv.org/abs/2606.11677v1. Full licence notice: \texttt{licenses/arxiv-2606.11677-CC-BY-4.0.txt}.\par\smallskip

\noindent\textit{Editorial scope.} Counted unit: the article's lemma lem:I1estimate (source0.tex lines 1897--1902). The article's complete original proof of it is delivered with it (source0.tex lines 1904--1961, one \texttt{proof} environment of 58 lines). No mathematics is written, repaired or re-derived: the statement and the proof are the article's own lines, in the article's own notation, and no retained line is substituted or re-flowed.  Retained uncounted as the source-local context the proof needs: lines 520--522; lines 1162--1165; lines 1171--1174; lines 1193--1198; lines 2352--2354.  Nothing was omitted from any retained span, so the package carries no omission record.  No retained line contains a citation, so the wrapper carries no bibliography and no bibliography entry.  No retained line contains a figure environment, an image-inclusion command or drawing code, so no image is delivered and the package declares no asset dependency.  Editorial formatting only, in addition: an AMS \texttt{amsart} preamble that restates the article's own declarations and macros used by the retained lines, namely the environments \texttt{lemma} on separate counters, together with the standard packages those lines need.  The retained environments print in this document as Lemma 1 in order of appearance, whereas the article prints the same items with the numbering of its own full document.  The complete native source of the article as distributed in the arXiv e-print for this exact version is bundled unchanged as \texttt{sources/202566/source0.tex}.
\begin{align}
H(a,x)\asymp \frac{C}{\max\{\sqrt{a},|x|\}}\label{eq:Hprop5}
\end{align}

By definition,  $\{M^\e_n\}_{n\geq 1}$  is a  reversible Markov chain with invariant measure and invariant distribution  \begin{align}
&\rho^\e(\dd x)=\rho^\e(x)\dd x=h^\e(x)^2V^\e(x)\dd x=p_{\e^2}(x)h^\e(x)\dd x=p_{\e^2}(x)\dd x\int_0^\infty\dd t \int_{\R^3}p_t(x,y)p_{\e^2}(y)\dd y\label{eq:invariantmeasure}\\
&\pi^\e(\dd x)=\pi^\e(x)\dd x:=2\pi^{\frac{3}{2}}\e \rho^\e(x)\dd x,\label{eq:invariantdistribution}
\end{align}

\begin{lemma}\label{lem:scalingMCM}
Suppose the Markov chain $\{M_n^1\}_{n\geq 0}$ starts at $M_0^1=x\in \R^3$. Then $\{\e M_n^1\}_{n\geq 0}$ is identically distributed as $\{M_n^\e\}_{n\geq 0}$ with $M_0^\e=\e x$.

\end{lemma}

\begin{lemma}\label{lem:MScorr}
Suppose  $M_0^{\e}=x\in \R^3$. Then, we have \begin{align*}
\mathtt{E}_x\left[f\left(M_n^\e\right)\right]=\mathbb{E}\left[f\left(S_n^{\e,x}\right)\right]
\end{align*}
for any $f\in B_b(\R^3)$ and $n\geq 1$. 
\end{lemma}

\begin{lemma}\label{lem:VuniGeoMe}
$\{M_n^1\}_{n\geq 0}$ is $V$-uniform integrable with $V(x)=\max\{|x|,1\}$.
\end{lemma}

\begin{lemma}\label{lem:I1estimate}
There exists a constant $C>0$ such that \begin{align*}
I_{\e,t,x}^{(1)}\leq C\e \frac{\delta}{a_\e}
\end{align*}
for uniformly in $|x|\geq r$.
\end{lemma}

\begin{proof}
First, we note that \begin{align*}
\lambda(\e)^\frac{s}{a_\e}=\left(1+\gamma{a_\e}\right)^\frac{t}{a_\e}\sim e^{\gamma t}
\end{align*}
as $\e\to 0$ for each $t\geq 0$. Also, for each $n\geq 0$,  
\begin{align*}
\lambda(\e)^n\leq e^{\frac{|\gamma|n}{a_\e}}.
\end{align*}

Thus, we obtain from Lemma \ref{lem:MScorr} that \begin{align*}
I_{\e,t,x}^{(1)}&\leq \sum_{n=1}^{\frac{\delta}{a_\e} }\lambda(\e)^nh^\e(x)\mathbb{E}\left[\frac{1}{h^\e(S_n^{\e,x})}\right]\leq Ch^\e(x)\sum_{n=1}^{\frac{\delta}{a_\e} }\mathtt{E}_x\left[\frac{1}{h^\e(M_n^{\e})}\right].
\end{align*}

Lemma \ref{lem:scalingMCM} and the Markovo proprty imply that 
\begin{align*}
\mathtt{E}_x\left[\frac{1}{h^\e(M_n^{\e})}\right]&=\e\mathtt{E}_{\frac{x}{\e}}\left[\frac{1}{h(M_n^{1})}\right]\\
&=\e \mathtt{E}_{\frac{x}{\e}}\left[\mathtt{E}_{M_1^1}\left[\frac{1}{h(M_{n-1}^{1})}\right]\right].
\end{align*}




By \eqref{eq:Hprop5}, we can apply Lemma \ref{lem:VuniGeoMe} and we can see that  there exists a constant $L>0$ such that \begin{align*}
\mathtt{E}_{M_1^1}\left[\frac{1}{h(M_{n-1}^{1})}\right]&= \int_{\R^3} \frac{1}{h(y)}\pi^1(\dd y)+R_n\\
&=2\pi^\frac{3}{2}+R_n,
\end{align*}
where where $\pi^1$ is given by \eqref{eq:invariantdistribution} and $R_n$ is a random variable with  
\begin{align*}
\left|R_n\right|\leq L\max\{|M_1^1|,1\} r^{n-1}.
\end{align*}
 

Also, we can see that for any $|x|\geq \e>0$ \begin{align*}
\mathtt{E}_{\frac{x}{\e}}\left[\max\{|M_1^1|,1\}\right]&\leq C\int_{\R^3}\frac{|y|p_1(y)}{\left|\frac{x}{\e}-y\right|h^1\left(\frac{x}{\e}\right)}+1\\
&\leq C\frac{|x|}{\e}\int_{\R^3}\frac{|y|p_1(y)}{\left|\frac{x}{\e}-y\right|}\dd y+1.
\end{align*}
Since we have \begin{align*}
\int_{\R^3}\frac{|y|p_1(y)}{\left|\frac{x}{\e}-y\right|}\dd y&=\int_{\left|\frac{x}{\e}-y\right|\leq \frac{|x|}{2\e}}\frac{|y|p_1(y)}{\left|\frac{x}{\e}-y\right|}\dd y+\int_{\left|\frac{x}{\e}-y\right|\geq \frac{|x|}{2\e}}\frac{|y|p_1(y)}{\left|\frac{x}{\e}-y\right|}\dd y\\
&\leq C\int_{\left|\frac{x}{\e}-y\right|\leq \frac{|x|}{2\e}}\frac{\frac{|x|}{2\e}e^{-\frac{|x|^2}{8\e^2}}}{\left|\frac{x}{\e}-y\right|}\dd y+\frac{2\e}{|x|}\int_{\R^3}|y|p_1(y)\dd y,
\end{align*}
it follows that \begin{align}
\mathtt{E}_{M_1^1}\left[\frac{1}{h(M_{n-1}^{1})}\right]=2\pi^\frac{3}{2}+r^{n-1}O(1).\label{eq:h(S)exp}
\end{align} 

Therefore, we can find a constant $C>0$ such that \begin{align*}
I_{\e,t,x}^{(1)}\leq C\e \frac{\delta}{a_\e}
\end{align*}
for uniformly in $|x|\geq r$.









\end{proof}

\end{document}
